% Este fichero es parte del Número 5 de la Revista Occam's Razor
% Revista Occam's Razor Número 5
%
% (c)  2010 The Occam's Razor Team
%
% Esta obra está bajo una licencia Reconocimiento 2.5 España de Creative
% Commons. Para ver una copia de esta licencia, visite 
% http://creativecommons.org/licenses/by/2.5/es/
% o envie una carta a Creative Commons, 171 Second Street, Suite 300, 
% San Francisco, California 94105, USA.

% Este .tex contiene el contenido del articulo
% Seccion (Dejar en blanco)
%

% Incluye imagen del artículo (Debe ser diferente del Kokopelli! (°_°) Sin embargo, no es obligatoria esta imagen)

\rput(2.5,-2.3){\resizebox{!}{3.2cm}{{\includegraphics[]{images/Importancia/glud}}}}

%Imagen de el comienzo de el articulo, coordenadas desde %la parte superior izquierda del margen de la pagina

% -------------------------------------------------
% Cabecera
\begin{flushright}
\msection{introcolor}{black}{0.35}{Seguridad Informática}

\mtitle{10cm}{Importancia del Software Libre en la Seguridad Informática}

\msubtitle{8cm}{Importance of Free Software in Computer Security.}

{\sf Diego Osorio, Andrés Acosta, Leidy Aldana. \\Estudiantes de pregrado Ingeniería de Sistemas}

{\psset{linecolor=black,linestyle=dotted}\psline(-12,0)}

\end{flushright}

\vspace{2mm}
% -------------------------------------------------

\begin{multicols}{2}

\sectiontext{white}{black}{Resumen}

% Introducción
\intro{introcolor}{R}{esumen pendiente!
%En esta parte por favor ubicar el Resumen de su artículo, debe permitir identificar en forma rápida el contenido básico del trabajo; no debe ser mayor a 250 palabras. No debe aportar información o conclusión que no está presente en el texto, así como tampoco debe  citar referencias bibliográficas. Debe quedar claro el problema que se investiga y el objetivo del mismo. 
}

\sectiontext{white}{black}{Palabras clave: } Confiabilidad, Integridad, Disponibilidad, Libertad

% Leidy
% Introducción, Importancia del Software Libre, Seguridad Informática
\vspace{1.0cm}

Compartir conocimiento ha aportado en la evolución de las sociedades, de forma similar, compartir código ha aportado a la construcción de software de mayor calidad; más allá del código abierto este artículo hace referencia al Software Libre, el cual permite la libertad del usuario y la libertad del software, permite modificar de acuerdo a necesidades propias, permite ayudar a una persona o a una comunidad; el software es libre al tener cuatro libertades, enunciadas en \textit{Software Libre para una Sociedad Libre} (Stallman, 2004, página 236): 

\vspace{0.1cm}

\begin{enumerate}
	\item[0.] La libertad cero es la libertad de ejecutar el programa con cualquier propósito. 
	\item[1.] La libertad uno es la libertad de ayudarte a ti mismo cambiando el programa para que se ajuste a tus necesidades. 
    \item[2.] La libertad dos es la libertad de ayudar al prójimo distribuyendo copias del programa. 
    \item[3.] La libertad tres es la de ayudar a distribuir tu comunidad publicando una versión mejorada de modo que los otros puedan beneficiarse de tu trabajo.
\end{enumerate}

\vspace{0.15cm}

Gracias a estas cuatro libertades el Software Libre ha contribuido y puede contribuir en el aprendizaje, en la construcción y el soporte de grandes empresas o en la enseñanza; de hecho puede ser una alternativa que resulta más interesante, alejada del tedio y más próxima al área científica. Cabe resaltar que libre es diferente de gratis (libre como pensamiento libre más no como barra libre), que el software libre no quita el reconocimiento de la persona que contribuye a crear software o otro aporte hecho a la construcción de conocimiento libre.

\vspace{0.3cm}

\begin{entradilla}
{\em El {\color{introcolor}{Software Libre}} tiene la ventaja de ser estudiado, ejecutado, modificado y distribuido.}
\end{entradilla}

\vspace{0.3cm}

Su importancia se ve reflejada en la mejora continua, en que brinda la oportunidad de aprender apartir de aportes realizados por otras personas, se ve reflejada en el hecho de que el conocimiento debería ser libre. Para algunas personas más allá de estas ventajas mencionadas anteriormente, la elección de utilizar Software Libre es una cuestión ética, de ciencia, de no permanecer tanto en la ignorancia, tomando como ejemplo el área de la seguridad informática (área que hace parte de la mayoría de las tecnologías de la información y la comunicación), hay programas no libres que cuentan con código oculto utilizado para conseguir información del usuario o ejecutar procesos sin permiso del usuario.\\

Se denomina por seguridad informática a las medidas y controles que garantizan la confiabilidad, integridad y disponibilidad de los activos del sistema de información; incluyendo hardware y software, firmware e información procesada, almacenada y comunicada. \\

El área de la seguridad informática influye en la mayoría de las áreas en la actualidad, como en el desarrollo web, en la inteligencia artificial, en el manejo de grandes volúmenes de información; desde una perspectiva más cotidiana influye en el simple hecho de enviar un correo electrónico (para que llegue sin modificaciones exclusivamente al destinatario), de hacer una transacción bancaria en linea, o realizar la conexión a una red wi-fi sin que la información de nuestros dispositivos sea vulnerable.

\vspace{0.3cm}

\begin{entradilla}
{\em La {\color{introcolor}{seguridad informática}} son las medidas y controles que garantizan la confiabilidad, integridad y disponibilidad de los activos del sistema de información.}
\end{entradilla}

De esta forma, este artículo busca mostrar al lector la importancia del trabajo conjunto de estos dos campos a través de ventajas, de herramientas, de alternativas.
\\ \\

% Diego -------------------------------------
% Implementación de herramientas
% confiabilidad, contribución, acceso, guias y documentación

Las herramientas utilizadas en la seguridad informática que son de código abierto, permiten que la comunidad examine el código, lo entiendan e incluso lo modifiquen, lo que permite a los usuarios aprender acerca del software, ajustarlo a sus necesidades e incluso contribuir. Esto lo hace más versátil e incluso agrega la confiabilidad de usarlo sabiendo que muchas personas han examinado el código y nos da la certeza que hace lo que debe, que en este no se encuentra código malicioso que pueda afectar a quien lo use y que ya ha sido probado.
\\
\\
El fácil acceso a todo tipo de usuarios, dónde no se tiene que pagar por utilizar, y el gran aumento del uso de las tecnologías de la información y el mundo digital, permite que la comunidad que utiliza software libre en la seguridad informática sea cada vez más, en los últimos años ha ganado mucha popularidad, incluso la demanda laboral ha aumentado drásticamente en los últimos años al igual que el uso de software libre a nivel profesional en la seguridad informática y personas especializadas en él.
\\
\\
Al ser probado constantemente por diferentes personas, con diferentes enfoques hace que se encuentren mayor número de problemas y fallos en menor tiempo que un software que sólo es mantenido por una empresa. Al haber más personas que pueden contribuir a mejorarlo, se ve una evolución constante haciendo las herramientas de software libre actualizadas, versátiles y altamente competitivas.
\\
\\
Es bastante la información que se encuentra en Internet de las herramientas de software libre en el campo de la seguridad informática, en los principales blogs se discute y comparte especialmente sobre software de código abierto. Tanto organizaciones como personas individuales comparten constantemente documentación, guías, tutoriales, vídeos, experiencias, preguntas, trucos, etc. generando una gran cantidad de material en múltiples idiomas permitiendo que casi cualquier persona con iniciativa y con algunos conocimientos básicos de la informática pueda aprender a utilizar estas herramientas.

%Andres

% Siguiente página
%%%%%%%%%%%%%%%%%%%%%%%%%%%%%%%%%%%%%%%%%%%%%%%%%%%%%%%%%%%%
%\ebOpage{introcolor}{0.35}{Seccion dejar vacio}
%%%%%%%%%%%%%%%%%%%%%%%%%%%%%%%%%%%%%%%%%%%%%%%%%%%%%%%%%%%%

% Dudas respecto al manejo del artículo
% https://github.com/GLUD/articulo-individual-2017/blob/master/nombre_articulo.tex

\medskip

\bibliographystyle{abbrv}
\begin{bibliografia}
\bibitem{kopka}
Stallman, Richard. \emph{Software Libre para una Sociedad Libre}, \href{https://www.gnu.org/philosophy/fsfs/free_software.es.pdf}{disponible en pdf}.\hskip 1em plus
  0.5em minus 0.4em\relax Madrid, España: Traficantes de sueños, 2004.

\end{bibliografia}

\begin{biografia}{images/Importancia/glud.eps}{Diego Osorio} 
% añadir fotografía tamaño [2.5 cm x 3.3 cm ] 
%  |¬_¬|'  La foto es Opcional! , se añade reemplazando:
% \begin{biografia}{images/mi_articulo/autor.eps}{Nombre Completo del Autor}
Corta descripción del autor, que incluya estudios de pregrado o postgrado, pertenencia a grupos de investigación
\end{biografia}

\begin{biografia}{images/Importancia/glud.eps}{Andrés Acosta} 
% añadir fotografía tamaño [2.5 cm x 3.3 cm ] 
%  |¬_¬|'  La foto es Opcional! , se añade reemplazando:
% \begin{biografia}{images/mi_articulo/autor.eps}{Nombre Completo del Autor}
Corta descripción del autor, que incluya estudios de pregrado o postgrado, pertenencia a grupos de investigación
\end{biografia}

\begin{biografia}{images/Importancia/white.eps}{Leidy Aldana} 
\\Estudiante de Ingeniería de Sistemas, interesada por aprender continuamente; quien ve en el software libre una perspectiva más agradable e interesante de su carrera.\\LeidyMarcelaAldana@gmail.com

\end{biografia}

\raggedcolumns
\pagebreak


\end{multicols}

\clearpage
\pagebreak
